% Preambulo compartido para clases en formato beamer
\documentclass[9pt, xcolor=svgnames]{beamer}
\usepackage[utf8]{inputenc}
\usepackage[T1]{fontenc}
\usepackage[spanish,es-nodecimaldot]{babel}

\usetheme[progressbar=frametitle]{metropolis}
\metroset{sectionpage=none, subsectionpage=none}
\usepackage{appendixnumberbeamer}
\usepackage{booktabs}
\usepackage{amsmath,amssymb,bm,mathtools}
\usepackage{physics}
\usepackage{graphicx}
\usepackage{multicol}
\usepackage{tikz}
\usepackage{minted}
\usemintedstyle{tango}
\setminted{fontsize=\scriptsize, breaklines=true, linenos=true, bgcolor=LightGray!20, frame=lines}
\newminted{python}{fontsize=\scriptsize, breaklines=true, linenos=true, bgcolor=LightGray!20, frame=lines}
\usepackage{pgfplots}
\pgfplotsset{compat=1.18}
\usepackage{ragged2e}
\usepackage{xspace}
\newcommand{\themename}{\textbf{\textsc{metropolis}}\xspace}

\setbeamercolor{block title}{bg=MidnightBlue!85, fg=white}
\setbeamercolor{block body}{bg=MidnightBlue!5, fg=black}
\setbeamercolor{alerted text}{fg=Crimson}
\setbeamercolor{frametitle}{bg=MidnightBlue!10, fg=black}
\setbeamertemplate{blocks}[rounded][shadow=true]


\weektitle{12}{redes informadas por fisica y cierre del curso}
\titlegraphic{}
\title[PCFI261 -- Semana 12]{\Large Modelos Computacionales PCFI261\\Semana 12: redes informadas por fisica\\y cierre del curso}

\newcommand{\compactcode}{\fontsize{6.0}{5.8}\selectfont}
\setminted[python]{fontsize=\compactcode,breaklines=true,breakanywhere=true,linenos=true,numbersep=3pt,bgcolor=LightGray!12,frame=single,framesep=1.1mm,rulecolor=MidnightBlue!55,baselinestretch=0.86,tabsize=2}

\begin{document}

\begin{frame}[plain]
\centering
\vfill
{\Large\bfseries Modelos Computacionales PCFI261}\\[0.25cm]
{\large Semana 12: redes informadas por fisica\\y cierre del curso}\\[0.7cm]
{\large Joaquin Peralta}\\[0.15cm]
{\small \texttt{joaquin.peralta@unab.cl}}\\[0.35cm]
{\small Universidad Andres Bello}
\vfill
\end{frame}

\begin{frame}{Organizacion de la semana}
\begin{block}{Clase final del curso}
La Semana 12 cierra el curso con un flujo completo: observacion visual, extraccion de datos, red neuronal con Keras y red informada por fisica. La pregunta central no es solo si una red ajusta datos, sino que significa que un modelo computacional respete una ley fisica.
\end{block}
\begin{columns}[T,onlytextwidth]
\column{0.48\textwidth}
\textbf{Clase 1}
\begin{itemize}
    \item de imagenes a trayectoria;
    \item datos con unidades y escala;
    \item red neuronal con \texttt{tensorflow.keras};
    \item curvatura aprendida y estimacion de $g$.
\end{itemize}
\column{0.48\textwidth}
\textbf{Clase 2}
\begin{itemize}
    \item perdida informada por fisica;
    \item derivadas con \texttt{tf.GradientTape};
    \item aprendizaje simultaneo de $y_\theta(t)$ y $g$;
    \item cierre: datos, modelo, validacion e interpretacion.
\end{itemize}
\end{columns}
\end{frame}

\begin{frame}{Objetivos de aprendizaje}
\begin{itemize}
    \item Reconstruir una trayectoria a partir de una secuencia simple de imagenes.
    \item Entrenar una red neuronal con \texttt{tensorflow.keras} para aproximar $x(t)$ e $y(t)$.
    \item Usar diferenciacion automatica de TensorFlow para calcular derivadas temporales.
    \item Construir una funcion de perdida que combine datos y una ley fisica.
    \item Comparar un modelo puramente flexible con un modelo informado por fisica.
    \item Sintetizar el ciclo completo de modelamiento computacional trabajado en el curso.
\end{itemize}
\end{frame}

\begin{frame}{Continuidad desde las semanas anteriores}
\begin{itemize}
    \item Semanas iniciales: discretizacion, integracion numerica, simulacion y errores.
    \item Semanas intermedias: modelos computacionales, analisis de datos y validacion.
    \item Semanas 09--11: redes neuronales, Keras, senales e imagenes.
    \item Semana 12: una red neuronal no reemplaza la fisica; la incorporamos como restriccion.
\end{itemize}
\begin{alertblock}{Idea de cierre}
Un modelo computacional cientifico no se evalua solo por ajustar datos. Tambien debe tener unidades, supuestos claros, validacion y una interpretacion fisica defendible.
\end{alertblock}
\end{frame}

\section{Clase 1}

\begin{frame}[plain]
\centering
\vfill
{\Huge \bfseries Clase 1}\\[0.4cm]
{\Large De imagenes a datos y de datos a una red neuronal}
\vfill
\end{frame}

\begin{frame}{Objetivos de la Clase 1}
\begin{itemize}
    \item Entender como se obtiene una tabla de trayectoria desde frames de video.
    \item Identificar los supuestos de escala, origen y segmentacion.
    \item Entrenar una red neuronal densa para aprender $x(t)$ e $y(t)$.
    \item Estimar la aceleracion vertical a partir de la curvatura aprendida por la red.
\end{itemize}
\end{frame}

\begin{frame}{Problema fisico y observacional}
\begin{columns}[T,onlytextwidth]
\column{0.52\textwidth}
\begin{block}{Situacion}
Observamos una pelota roja en movimiento. No partimos desde una tabla ideal: primero debemos convertir imagenes en datos numericos.
\end{block}
\begin{itemize}
    \item Entrada inicial: frames del movimiento.
    \item Dato intermedio: centro de la pelota en pixeles.
    \item Dato fisico: posicion en metros y tiempo en segundos.
    \item Modelo esperado: movimiento de proyectil sin roce.
\end{itemize}
\column{0.44\textwidth}
\begin{alertblock}{Supuestos}
\begin{itemize}
    \item camara fija;
    \item escala conocida;
    \item objeto segmentable por color;
    \item perspectiva despreciable.
\end{itemize}
\end{alertblock}
\end{columns}
\end{frame}

\begin{frame}[fragile]{Notebook 01: extraer trayectoria}
\begin{minted}{python}
mask = (arr[:, :, 0] > 180) & (arr[:, :, 1] < 120) & (arr[:, :, 2] < 120)
ys, xs = np.where(mask)
cx = xs.mean(); cy = ys.mean()
x_m = (cx - origin_px[0]) / px_per_m
y_m = (origin_px[1] - cy) / px_per_m
\end{minted}
\begin{block}{Lectura docente}
La parte importante no es que el detector sea sofisticado. La parte importante es reconocer donde entran los supuestos que convierten pixeles en magnitudes fisicas.
\end{block}
\end{frame}

\begin{frame}{De datos a modelo flexible}
\begin{itemize}
    \item Con la tabla \texttt{t, x, y}, podemos ajustar un modelo.
    \item Un ajuste polinomial seria natural para un proyectil ideal.
    \item Esta semana preguntamos que aprende una red si no le damos la forma de la parabola.
    \item Usamos una red densa pequena con entrada $t$ y salida $(x,y)$.
\end{itemize}
\begin{alertblock}{Decision de implementacion}
Se usa \texttt{tensorflow.keras} para mantener continuidad con las semanas previas del curso y evitar cambios de backend innecesarios.
\end{alertblock}
\end{frame}

\begin{frame}[fragile]{Notebook 02: red neuronal con Keras}
\begin{minted}{python}
from tensorflow import keras
from tensorflow.keras import layers
model = keras.Sequential([
    layers.Input(shape=(1,)),
    layers.Dense(32, activation="tanh"),
    layers.Dense(32, activation="tanh"),
    layers.Dense(2)
])
model.compile(optimizer=keras.optimizers.Adam(0.01), loss="mse")
\end{minted}
\end{frame}

\begin{frame}{Que esta aprendiendo la red}
\begin{columns}[T,onlytextwidth]
\column{0.50\textwidth}
\begin{itemize}
    \item La red aproxima una funcion continua.
    \item La entrada es un escalar: tiempo.
    \item La salida tiene dos componentes: posicion horizontal y vertical.
    \item El ajuste se evalua contra datos, no contra una ecuacion fisica.
\end{itemize}
\column{0.46\textwidth}
\begin{block}{Pregunta clave}
Si el error contra datos es pequeno, podemos afirmar que la red aprendio la ley del movimiento?
\end{block}
\begin{alertblock}{Respuesta corta}
No necesariamente. Ajustar puntos e identificar una ley son tareas relacionadas, pero no equivalentes.
\end{alertblock}
\end{columns}
\end{frame}

\begin{frame}[fragile]{Curvatura mediante TensorFlow}
\begin{minted}{python}
t_tf = tf.convert_to_tensor(t.reshape(-1, 1), dtype=tf.float32)
with tf.GradientTape() as tape2:
    tape2.watch(t_tf)
    with tf.GradientTape() as tape1:
        tape1.watch(t_tf)
        y_out = model(t_tf, training=False)[:, 1:2]
    dy_dt = tape1.gradient(y_out, t_tf)
d2y_dt2 = tape2.gradient(dy_dt, t_tf)
g_est = -float(tf.reduce_mean(d2y_dt2).numpy())
\end{minted}
\begin{block}{Interpretacion}
La diferenciacion automatica permite inspeccionar la curvatura de la funcion aprendida. Si $d^2y/dt^2 \approx -g$, la red es compatible con un tiro ideal.
\end{block}
\end{frame}

\begin{frame}{Actividad guiada de la Clase 1}
\begin{enumerate}
    \item Ejecutar \texttt{01\_extraccion\_trayectoria.ipynb}.
    \item Revisar la tabla \texttt{trajectory\_extracted.csv}.
    \item Ejecutar \texttt{02\_aprendizaje\_con\_red\_neuronal.ipynb}.
    \item Comparar datos y prediccion de la red.
    \item Estimar $g$ desde la segunda derivada aprendida.
    \item Discutir si el resultado es estable al cambiar epocas, neuronas o tasa de aprendizaje.
\end{enumerate}
\end{frame}

\begin{frame}{Cierre Clase 1}
\begin{itemize}
    \item Ya recorrimos imagenes $\rightarrow$ datos $\rightarrow$ red neuronal.
    \item La red puede interpolar una trayectoria simple con bajo error.
    \item La curvatura aprendida permite hacer una primera lectura fisica.
    \item Falta una pieza: obligar al modelo a respetar la ecuacion durante el entrenamiento.
\end{itemize}
\end{frame}

\section{Clase 2}

\begin{frame}[plain]
\centering
\vfill
{\Huge \bfseries Clase 2}\\[0.4cm]
{\Large Red informada por fisica y cierre del curso}
\vfill
\end{frame}

\begin{frame}{Objetivos de la Clase 2}
\begin{itemize}
    \item Definir una red $y_\theta(t)$ para la coordenada vertical.
    \item Introducir un parametro aprendible $g$.
    \item Construir una perdida con datos y residuo fisico.
    \item Entrenar el modelo con un bucle explicito de TensorFlow/Keras.
    \item Discutir interpretabilidad, extrapolacion y limites del enfoque.
\end{itemize}
\end{frame}

\begin{frame}{De red libre a red informada}
\begin{columns}[T,onlytextwidth]
\column{0.48\textwidth}
\textbf{Red libre}
\begin{itemize}
    \item minimiza error contra datos;
    \item puede interpolar bien;
    \item no sabe que existe gravedad;
    \item su curvatura se revisa despues.
\end{itemize}
\column{0.48\textwidth}
\textbf{Red informada por fisica}
\begin{itemize}
    \item minimiza error contra datos;
    \item penaliza violar la ecuacion;
    \item aprende un parametro fisico;
    \item usa la ley durante el entrenamiento.
\end{itemize}
\end{columns}
\end{frame}

\begin{frame}{Perdida informada por fisica}
Para el movimiento vertical ideal:
\[
    \frac{d^2y}{dt^2} = -g.
\]
Definimos una red $y_\theta(t)$ y un parametro aprendible $g$. La perdida total es:
\[
    \mathcal{L}=\frac{1}{N}\sum_i \left(y_\theta(t_i)-y_i\right)^2+
    \alpha\frac{1}{N}\sum_i \left(\frac{d^2y_\theta}{dt^2}(t_i)+g\right)^2.
\]
\begin{block}{Rol de $\alpha$}
Controla cuanto confiamos en la ley fisica respecto de los datos observados. En un laboratorio real, elegir $\alpha$ tambien es una decision de modelamiento.
\end{block}
\end{frame}

\begin{frame}[fragile]{Notebook 03: parametro fisico aprendible}
\begin{minted}{python}
model = keras.Sequential([
    layers.Input(shape=(1,)),
    layers.Dense(32, activation="tanh"),
    layers.Dense(32, activation="tanh"),
    layers.Dense(1)
])
g_param = tf.Variable([9.0], dtype=tf.float32, name="g_param")
optimizer = keras.optimizers.Adam(learning_rate=0.01)
\end{minted}
\begin{alertblock}{Nota}
El parametro $g$ se actualiza junto con los pesos de la red. No es un valor fijo impuesto a mano.
\end{alertblock}
\end{frame}

\begin{frame}[fragile]{Derivadas para el residuo fisico}
\begin{minted}{python}
def segunda_derivada_modelo(t_in, training=True):
    with tf.GradientTape() as tape2:
        tape2.watch(t_in)
        with tf.GradientTape() as tape1:
            tape1.watch(t_in)
            y_out = model(t_in, training=training)
        dy_dt = tape1.gradient(y_out, t_in)
    d2y_dt2 = tape2.gradient(dy_dt, t_in)
    return y_out, d2y_dt2
\end{minted}
\end{frame}

\begin{frame}[fragile]{Bucle de entrenamiento}
\begin{minted}{python}
with tf.GradientTape() as tape:
    y_pred = model(t_tensor, training=True)
    loss_data = tf.reduce_mean(tf.square(y_pred - y_tensor))
    _, d2y_dt2 = segunda_derivada_modelo(t_tensor, training=True)
    physics_res = d2y_dt2 + g_param
    loss_phys = tf.reduce_mean(tf.square(physics_res))
    loss = loss_data + alpha * loss_phys
variables = model.trainable_variables + [g_param]
grads = tape.gradient(loss, variables)
optimizer.apply_gradients(zip(grads, variables))
\end{minted}
\end{frame}

\begin{frame}{Como leer el resultado}
\begin{itemize}
    \item Si el error de datos baja, la red esta ajustando las observaciones.
    \item Si el residuo fisico baja, la red satisface mejor la ecuacion.
    \item Si $g$ converge cerca de $9.8~\mathrm{m/s^2}$, el parametro aprendido es fisicamente razonable.
    \item Si el ajuste visual es bueno pero el residuo es grande, el modelo interpola sin respetar la ley.
\end{itemize}
\begin{alertblock}{Cuidado}
Una PINN no vuelve verdaderos los supuestos. Si hay roce, mala escala o mala deteccion, la ecuacion ideal puede ser una restriccion equivocada.
\end{alertblock}
\end{frame}

\begin{frame}{Actividad guiada de la Clase 2}
\begin{enumerate}
    \item Ejecutar \texttt{03\_red\_informada\_por\_fisica.ipynb}.
    \item Registrar el valor final aprendido de $g$.
    \item Comparar grafico de $y(t)$ y residuo fisico.
    \item Cambiar $\alpha$ y observar el compromiso entre ajuste y ley.
    \item Comparar con la red libre del notebook 02.
    \item Discutir que pasaria con datos mas ruidosos o con resistencia del aire.
\end{enumerate}
\end{frame}

\begin{frame}{Sintesis del curso}
\begin{columns}[T,onlytextwidth]
\column{0.50\textwidth}
\textbf{Habilidades computacionales}
\begin{itemize}
    \item programar modelos reproducibles;
    \item discretizar ecuaciones;
    \item simular y analizar resultados;
    \item entrenar modelos con Keras;
    \item validar con metricas y graficos.
\end{itemize}
\column{0.46\textwidth}
\textbf{Habilidades cientificas}
\begin{itemize}
    \item explicitar supuestos;
    \item revisar unidades y escalas;
    \item distinguir ajuste de explicacion;
    \item interpretar parametros;
    \item comunicar limitaciones.
\end{itemize}
\end{columns}
\end{frame}

\begin{frame}{Cierre de PCFI261}
\begin{block}{Mensaje final}
Modelar computacionalmente no es solo escribir codigo. Es construir una cadena razonada entre un fenomeno, una representacion numerica, un algoritmo, una validacion y una interpretacion.
\end{block}
\begin{itemize}
    \item La simulacion permite explorar consecuencias de leyes conocidas.
    \item El aprendizaje automatico permite ajustar patrones desde datos.
    \item La ciencia computacional aparece cuando unimos ambas cosas con criterio fisico.
    \item Esa es la idea central que esta semana deja instalada para proyectos futuros.
\end{itemize}
\end{frame}

\end{document}
